
This thesis presents \emph{thelab-mcp}, a local-first, multi-agent orchestration architecture built on the Model Context Protocol (MCP), a standardized communication protocol between LLM-based agents and local or external data environments. The architecture presented is an autonomous machine learning (ML) factory that converts a local tabular dataset into a versioned, validated, locally served ML model, with a single human approval gate and no mandatory cloud dependency.

The system is designed around the separation of responsibilities. Large language models decide \emph{what} to run, while a deterministic Python pipeline executes \emph{how} it runs. Typed contracts (Pydantic), a SQLite+FTS5 context store, seven MCP servers, an AST-restricted code sandbox, and a single human approval gate connect the agentic layer to the deterministic core. Agents (an analyst, a feature engineer, and a model selector, coordinated through an orchestrator) can only propose and compute inside bounded, sandboxed rounds. Every execution is recorded as evidence in run manifests, every claim an agent makes must resolve to recorded metrics, and a rejected proposal can never be executed. The design principle is \emph{grounded autonomy}. The deterministic factory provides the grounding evidence, and the agents operate on top of it as an additional layer.

The evaluation answers six research questions in two groups. The \emph{system group} is verified by an automated evaluator that runs inside a disposable temporary workspace. (RQ1) A model-training run is reproducible from its dataset, configuration, and seed. (RQ2) An independent MCP client can discover and use a registered model without any knowledge of the training internals. (RQ3) Local context retrieval recovers past evidence without modifying the store. The \emph{agentic group} measures what the agent layer adds on top of that grounding. (RQ4) Context-grounded agents produce more verifiable proposals than ungrounded agents. (RQ5) Sandboxed agent-generated code trains valid, competitive models beyond the fixed deterministic grid. (RQ6) Role-specialized agents coordinated via MCP outperform a single shared-prompt agent.

Different experiments complemented the automated checks. A three-ways-in comparison (CLI, HTTP, MCP) showed identical results across the interfaces. A real-dataset validation demonstrated correct and traceable exploratory data analysis and cleaning policies. Grounded and ungrounded agentic rounds showed that context produces actual verifiable claims instead of guesses. Sandboxed agent code managed to train valid and competitive models beyond the fixed grid. Role-specialized coordination outperformed single shared prompts. Multi-domain stress validation across ten datasets, from genomics and health to crypto, energy, and finance, surpassed the baseline models and generated an absorption loop to improve the deterministic pipeline. A fully-offline execution with a local 3B model demonstrated the architecture's independence from any cloud provider.

The results support the hypothesis. The architecture executes reproducible data-to-model workflows and exposes their artifacts to independent agents without being limited to any specific LLM provider. Its limitations (sandbox filesystem confinement, stdio-only transport, dataset-specific agentic findings) are documented as concrete, traceable findings rather than omitted.

\textbf{Keywords:} Multi-agent systems, Model Context Protocol, reproducibility, machine learning pipelines, interoperability, context engineering, local-first software.
